% answers_short.tex —— 2002 年·全国旧课程卷·文科 · 简短答案速查
% 由 tools/make_short_answers.py 自 answers.tex 的 \gksolhead/\gkanswer 提取，
% 并与 2002/2002_数学真题_全国旧课程文_答案要点.md 逐题比对后生成，勿手改（改源 MD 后重跑）。

\begin{gknotice}
\textbf{使用说明}：本节只给最终答案便于快速核对，完整推导见解析版。答案为编写者独立推导，\textbf{非官方标准答案}。
\end{gknotice}

\gkpart{选择题}
\begin{center}\begin{tabular}{ccccccccccccc}
\toprule
题号 & 1 & 2 & 3 & 4 & 5 & 6 & 7 & 8 & 9 & 10 & 11 & 12 \\
\midrule
答案 & D & C & D & B & C & B & B & C & D & A & D & B \\
\bottomrule
\end{tabular}\end{center}

\gkpart{填空题}
\begin{itemize}[leftmargin=2.6em]
  \item[\textbf{13.}] 1995—2000
  \item[\textbf{14.}] $(0,0),(1,1)$
  \item[\textbf{15.}] $1008$
  \item[\textbf{16.}] ②⑤
\end{itemize}

\gkpart{解答题}
\begin{itemize}[leftmargin=2.6em]
  \item[\textbf{17.}] (1) $20^\circ\mathrm C$；(2) $T=10\sin\left(\dfrac\pi8(t-10)\right)+20$，$6\leqslant t\leqslant14$
  \item[\textbf{18.}] (1) 7 分钟；(2) 15 分钟
  \item[\textbf{19.}] (1) $\dfrac{\sqrt3}{3}a^3$；(2) 见解析版
  \item[\textbf{20.}] (1) 既不是奇函数，也不是偶函数；(2) $\dfrac34$
  \item[\textbf{21.}] $y=x-1$ 或 $y=-x+1$。
  \item[\textbf{22.}] (1) 见解析版中的一组剪拼方案；(2) 正三棱柱的体积大于正三棱锥的体积；(3) 见解析版中的一组剪拼方案
\end{itemize}
